\begin{pastexam}[label={px:ncku-113-prog-15}]{成大資訊聯招 113 學年度 第 15 題}
（計算）

(10\%) Find the set of feasible solutions $\{x=(x_1,x_2,x_3,x_4,x_5)\}$ or determine that no
feasible solution exists for the following system of difference constraints:
\begin{align*}
x_1-x_3&\le 1\\
x_2-x_3&\le 4\\
x_4-x_5&\le -2\\
x_3-x_4&\le 4\\
x_5-x_1&\le 3\\
x_4-x_2&\le -7\\
x_1-x_2&\le -2\\
x_5-x_3&\le 1
\end{align*}
\end{pastexam}

\begin{pastsoln}
答案：有可行解，例如 $x = (-2, 0, -3, -7, -2)$。

題目寫「the set of feasible solutions」。 一個差分約束系統有解時，解有無限多組（資結書引理~25.8：
任何一組解每個分量都加同一個數仍是解），而且解集合不只這些平移，所以這裡的解讀是「求出一組可
行解」，再說明怎麼得到更多組。

\textbf{約束圖}（資結書第~25.4~節）。 頂點 $v_0, v_1, \ldots, v_5$；$v_0$ 到每個 $v_i$ 一條權
重 $0$ 的邊；每條 $x_j - x_i \le b$ 是邊 $(v_i, v_j)$，權重 $b$。 依題目的順序：
\begin{center}
\small
\begin{tabular}{llll}
\toprule
限制 & 邊 & 限制 & 邊 \\
\midrule
$x_1 - x_3 \le 1$ & $(v_3, v_1)$，$1$ & $x_5 - x_1 \le 3$ & $(v_1, v_5)$，$3$ \\
$x_2 - x_3 \le 4$ & $(v_3, v_2)$，$4$ & $x_4 - x_2 \le -7$ & $(v_2, v_4)$，$-7$ \\
$x_4 - x_5 \le -2$ & $(v_5, v_4)$，$-2$ & $x_1 - x_2 \le -2$ & $(v_2, v_1)$，$-2$ \\
$x_3 - x_4 \le 4$ & $(v_4, v_3)$，$4$ & $x_5 - x_3 \le 1$ & $(v_3, v_5)$，$1$ \\
\bottomrule
\end{tabular}
\end{center}
（左欄由上而下是第 1--4 條，右欄是第 5--8 條。）

\textbf{Bellman--Ford}，從 $v_0$ 出發，每一輪先掃 $v_0$ 的五條邊，再依上表的順序掃：
\begin{center}
\small
\begin{tabular}{c*{6}{c}l}
\toprule
輪 & $d[v_0]$ & $d[v_1]$ & $d[v_2]$ & $d[v_3]$ & $d[v_4]$ & $d[v_5]$ & 這一輪成功鬆弛的限制邊 \\
\midrule
1 & $0$ & $-2$ & $0$ & $0$ & $-7$ & $0$ & $(v_5, v_4)$、$(v_2, v_4)$、$(v_2, v_1)$ \\
2 & $0$ & $-2$ & $0$ & \underline{$-3$} & $-7$ & \underline{$-2$} & $(v_4, v_3)$、$(v_3, v_5)$ \\
3 & $0$ & $-2$ & $0$ & $-3$ & $-7$ & $-2$ & 沒有 \\
\bottomrule
\end{tabular}
\end{center}
（有底線的是第 2 輪變小的值。） 第 1 輪掃完 $v_0$ 的邊後 $d[v_1..v_5]$ 都是 $0$；$(v_5, v_4)$ 先把 $d[v_4]$ 改成 $-2$，$(v_2, v_4)$ 再
改成 $-7$，$(v_2, v_1)$ 把 $d[v_1]$ 改成 $-2$；$(v_3, v_1)$ 在它們之前掃到，$0 + 1$ 不比 $0$ 小。 第 2 輪
$(v_4, v_3)$：$-7 + 4 = -3$；$(v_3, v_5)$ 用剛改好的 $d[v_3]$：$-3 + 1 = -2$。 第 3 輪沒有改變，之後
第 4、5 輪和檢查也不會有，回傳 \TRUE：約束圖沒有負環，系統有解。

\textbf{一組解與驗算。} $x_i = d[v_i]$，$x = (-2, 0, -3, -7, -2)$。 代回八條限制：
$x_1 - x_3 = 1 \le 1$；$x_2 - x_3 = 3 \le 4$；$x_4 - x_5 = -5 \le -2$；$x_3 - x_4 = 4 \le 4$；
$x_5 - x_1 = 0 \le 3$；$x_4 - x_2 = -7 \le -7$；$x_1 - x_2 = -2 \le -2$；$x_5 - x_3 = 1 \le 1$，全部成
立。

\textbf{更多的解。} 對任何實數 $c$，$(-2 + c,\ c,\ -3 + c,\ -7 + c,\ -2 + c)$ 都是解，例如 $c = 7$
得到全部非負的 $(5, 7, 4, 0, 5)$。 完整的解集合就是滿足這八條不等式的所有 $x$；在 $x_i \le 0$ 的解裡，
上面這組每個分量都最大（資結書第~25.4~節）。
\end{pastsoln}
